\documentclass[11pt]{article}
\usepackage[a4paper,margin=27mm]{geometry}
\usepackage{amsmath,amssymb,amsthm,booktabs,array,longtable}
\usepackage[T1]{fontenc}
\usepackage{lmodern}
\usepackage[colorlinks=true,allcolors=blue,breaklinks=true]{hyperref}
\setlength{\emergencystretch}{1em}
\newtheorem{theorem}{Theorem}
\newtheorem{lemma}[theorem]{Lemma}
\newtheorem{corollary}[theorem]{Corollary}
\newcommand{\R}{\mathbb R}
\newcommand{\N}{\mathbb N}
\newcommand{\diag}{\operatorname{diag}}
\newcommand{\tr}{\operatorname{tr}}
\newcommand{\supp}{\operatorname{supp}}
\title{Koljan\v{c}i\'c--Oblak Question 5.11\\
(arXiv:2604.21552v1, Section 5):\\
nSSP witnesses for looped double paths}
\author{Haobo Ma\\\small ChronoAI Pte Ltd, Singapore
\and Wenlin Zhang\\\small National University of Singapore\\
\small ChronoAI Pte Ltd, Singapore}
\date{November 2, 2026}
\begin{document}
\maketitle
\begin{abstract}
Question 5.11 of Koljan\v{c}i\'c and Oblak
(arXiv:2604.21552v1, Section 5; journal p.~776) has an affirmative answer:
every even double path
with a nonempty loop set allows the non-symmetric Strong Spectral
Property. We construct a symmetric one-loop witness using coprime
branch continuants and prove injectivity of its characteristic-polynomial
differential. Openness then supplies every prescribed nonempty loop set.
The same argument gives a complete existence criterion in odd order:
at least one loop must lie at an odd-numbered vertex. The witnesses can
be chosen symmetric with positive integer nonzero entries. Carptopus
(19 August 2026) previously proved the full existence criterion and its
single-loop spectral characterization. We give an explicit integer
construction, an arithmetic criterion for unit edge weights, and the
three-way classification for the central-loop and interior-loop families.
\end{abstract}

\section{The question and the result}
Write $[n]=\{1,\ldots,n\}$, with $n\geq1$. The double path
$P_{n,L}$ has both arcs between consecutive vertices and exactly the
loops at $L\subseteq[n]$. Following Definitions 1.1, 1.3 and 1.5 of
\cite{KO}, a real matrix $A$ has nSSP when
\begin{equation}\label{eq:test}
 A\circ X=0,\qquad AX^{\mathsf T}=X^{\mathsf T}A
 \quad\Longrightarrow\quad X=0\qquad(X\in\R^{n\times n}).
\end{equation}
The product $\circ$ is entrywise. A pattern \emph{allows} nSSP if some
matrix with exactly that support has nSSP; it \emph{requires} nSSP if
every such matrix does. No symmetry or diagonal condition is imposed on
the test matrix $X$.

Question 5.11 of \cite{KO} is the even-order existence question.
Its answer follows from the following stronger statement.
\begin{theorem}\label{thm:main}
The double path $P_{n,L}$ allows nSSP if and only if
\[
 \begin{cases}
 L\ne\varnothing,&n\text{ even},\\
 L\cap\{1,3,\ldots,n\}\ne\varnothing,&n\text{ odd}.
 \end{cases}
\]
Whenever it allows nSSP, it has a symmetric nSSP witness whose nonzero
entries are positive integers.
\end{theorem}
Our focus is existence, not the universal classification in Problem 1.6
of \cite{KO}. In particular, the negative requiring examples there
are compatible with Theorem~\ref{thm:main}. The odd-order obstruction
is already Theorem 5.7(Cc) of that paper; we include its short linear
algebra proof to make the criterion self-contained.
Carptopus \cite[Theorems 4.1, 5.2 and 6.1]{C} gave an earlier proof of
both parity cases, dated 19 August 2026 and revised 22 August 2026.
The sufficient direction uses the same separation of continuant
variations into their two parities. The construction below fixes the
branch scaling explicitly and retains the positive-integer conclusion
after every loop enlargement. For context, disjoint branch spectra also
appear in the three-spectra Jacobi reconstruction theorem of
Michor and Teschl \cite[Theorem 2.1]{MT}; here the argument establishes
the required differential directly.
Vander Meulen's ILAS 2026 abstract, reporting joint work with Cavers
and Li \cite[p.~71]{VM}, announces a general sign-pattern
characterization; the abstract does not give its hypotheses or proof.

\section{Continuants and their differential}
For positive squared edge weights $q_0,q_1,\ldots$, put
\begin{equation}\label{eq:rec}
 P_0=1,\quad P_1=x,\quad
 P_{k+2}=xP_{k+1}-q_kP_k.
\end{equation}
For a branch of $m$ vertices, the first $m-1$ weights describe its
internal edges; $q_{m-1}$ describes the edge attaching it to a central
vertex. The polynomial $P_m$ is monic of degree $m$ and has parity $m$.
For a variation of the $q_k$, denote differentiation by a dot. Then
$\dot P_0=\dot P_1=0$ and
\begin{equation}\label{eq:drec}
 \dot P_{k+2}=x\dot P_{k+1}-\dot q_kP_k-q_k\dot P_k.
\end{equation}
Thus $\deg\dot P_m\leq m-2$ for $m\geq2$, and $\dot P_m$ has parity
$m$; zero polynomials satisfy all the degree bounds below.
These assertions follow directly by induction in \eqref{eq:rec}
and \eqref{eq:drec}.

\begin{lemma}\label{lem:recover}
If $q_0,\ldots,q_{m-1}\ne0$ and
$\dot P_m=\dot P_{m+1}=0$, then every
$\dot q_0,\ldots,\dot q_{m-1}$ vanishes.
\end{lemma}
\begin{proof}
For $m=0$ there is nothing to prove. For $m\geq1$, the recurrence gives
$\dot q_{m-1}P_{m-1}+q_{m-1}\dot P_{m-1}=0$.
Its coefficient of $x^{m-1}$ is $\dot q_{m-1}$, by monicity and the
strict degree bound. Hence $\dot q_{m-1}=0$ and, since
$q_{m-1}\ne0$, also $\dot P_{m-1}=0$. Descend in $m$.
\end{proof}

Choose a loop vertex $\ell$, and put $r=\ell-1$, $s=n-\ell$.
Use two copies $P_k,Q_k$ of \eqref{eq:rec}, reading each branch from
its far endpoint toward the loop. Let $U=P_r$, $V=Q_s$.
The symmetric path matrix with loop value $a\ne0$ and edge weights
equal to the positive square roots of the parameters has characteristic
polynomial
\begin{equation}\label{eq:joined}
 H=P_{r+1}V+UQ_{s+1}-(x+a)UV.
\end{equation}
Indeed, expansion at the central vertex gives
$(x-a)UV-q_{r-1}P_{r-1}V-p_{s-1}UQ_{s-1}$, with a branch term
omitted when its length is zero. Substitution of \eqref{eq:rec}
gives \eqref{eq:joined}; in particular the formula includes empty branches.

\begin{lemma}\label{lem:joined}
If $a\ne0$, $U,V$ are coprime and all the edge parameters are nonzero, then
$\dot H=0$ forces $\dot a=0$ and every edge-parameter variation to vanish.
\end{lemma}
\begin{proof}
Set $T=\dot U V+U\dot V$. Differentiating \eqref{eq:joined} gives
\[
 \dot H=W-aT-\dot aUV,
\]
where
\[
 W=\dot P_{r+1}V+P_{r+1}\dot V+
       \dot UQ_{s+1}+U\dot Q_{s+1}-xT.
\]
Every term of $W$ has parity $r+s+1$, whereas $T$ and $UV$ have
parity $r+s$. The even and odd parts of $\dot H=0$ therefore give
$aT+\dot aUV=0$ and $W=0$. The coefficient of $x^{r+s}$ in the first
identity gives $\dot a=0$, because $UV$ is monic and
$\deg T<r+s$. Since $a\ne0$, $T=0$.

Coprimality now gives $U\mid\dot U$ and $V\mid\dot V$.
The bounds $\deg\dot U<\deg U$ and
$\deg\dot V<\deg V$ force $\dot U=\dot V=0$.
Consequently $W=0$ reduces to
\[
 \dot P_{r+1}V+U\dot Q_{s+1}=0.
\]
The same divisibility argument, using
$\deg\dot P_{r+1}<r$ and $\deg\dot Q_{s+1}<s$, gives
$\dot P_{r+1}=\dot Q_{s+1}=0$.
For a zero-length branch these variations are identically zero, so
the argument still applies with its constant polynomial $1$.
Lemma~\ref{lem:recover} finishes both branches.
\end{proof}
This parity argument applies whenever the branches are coprime.

\section{Choosing the branches and proving nSSP}
\begin{lemma}\label{lem:choose}
If at least one of $r,s$ is even, the branches can be chosen with
positive integer edge weights and coprime characteristic polynomials.
\end{lemma}
\begin{proof}
An empty branch has polynomial $1$, so suppose both are nonempty.
Choose an even branch of length $d\geq2$. Give all its internal edges,
and its attaching edge, weight $M=d+1$; give the other branch all weights
$1$. The $d$ eigenvalues of a zero-diagonal path with constant weight
$M$ are
$2M\cos(j\pi/(d+1))$, $1\leq j\leq d$.
This formula follows by substituting the vectors with coordinates
$\sin(ij\pi/(d+1))$ in the three-term matrix equation; the displayed
eigenvalues are distinct and hence exhaust its spectrum.
For even $d$, their least absolute value is
\[
 2(d+1)\sin\!\left(\frac{\pi}{2(d+1)}\right)>2,
\]
using $\sin t>2t/\pi$ for $0<t<\pi/2$.
Every eigenvalue of the other, unit-weight branch has absolute value
less than $2$ by the same formula. The spectra are disjoint. Since
both characteristic polynomials split over $\R$, they are coprime.
\end{proof}

Let $A$ be the resulting symmetric one-loop matrix, with $a=1$.
All its adjacent entries are nonzero. Such a tridiagonal matrix has
simple spectrum: its first coordinate determines an eigenvector
recursively from the first row onward, and a zero first coordinate
forces the entire vector to vanish. Every eigenspace therefore has
dimension at most one; the real symmetric spectral theorem then gives
an orthogonal matrix $O$ with
$O^{\mathsf T}AO=\diag(\lambda_1,\ldots,\lambda_n)$ and distinct
$\lambda_i$.

Let $\mathcal S$ be the $n$-dimensional real space of symmetric
matrices supported on the $n-1$ path edges and the chosen loop.
Consider
\begin{equation}\label{eq:readout}
 B\longmapsto\diag(O^{\mathsf T}BO)\in\R^n,
 \qquad B\in\mathcal S.
\end{equation}
This map is injective. To see this, suppose its output is zero.
For every $x$, the derivative at $t=0$ of
$\det(xI-A-tB)$ is
\[
 -\sum_i (O^{\mathsf T}BO)_{ii}\prod_{j\ne i}(x-\lambda_j)=0;
\]
the formula follows by multilinearity of the determinant at a diagonal
matrix. Thus $\dot H=0$. An edge variation $\dot b$ gives the squared
parameter variation $\dot q=2b\dot b$; since $b\ne0$,
Lemma~\ref{lem:joined} forces $B=0$. Equal dimensions make
\eqref{eq:readout} surjective as well.

Now let an unrestricted real $X$ satisfy \eqref{eq:test}.
The matrix $Y=O^{\mathsf T}X^{\mathsf T}O$ commutes with
$\diag(\lambda_i)$, so its $(i,j)$ entry is zero for $i\ne j$.
Write $Y=\diag(y_1,\ldots,y_n)$.
For every $B\in\mathcal S$, entrywise annihilation gives
$\tr(X^{\mathsf T}B)=\sum_{ij}X_{ij}B_{ij}=0$.
Choose $B$ whose output in \eqref{eq:readout} is $(y_i)$.
Invariance of trace under conjugation yields
\[
 0=\tr\bigl(YO^{\mathsf T}BO\bigr)=\sum_i y_i^2.
\]
Hence $Y=0$ and $X=0$. This proves nSSP for the one-loop witness
without assuming symmetry of $X$.

\begin{corollary}\label{cor:spectral}
For a real symmetric irreducible tridiagonal matrix with exactly one
nonzero diagonal entry, at $\ell$, nSSP holds if and only if the
two branches obtained by deleting $\ell$ have disjoint spectra.
In particular, if every edge has weight $1$ and the loop has any
nonzero value, then nSSP holds exactly when $\gcd(\ell,n+1)=1$.
\end{corollary}
\begin{proof}
The sufficient direction was proved above and is unchanged for signed
nonzero edges, whose squares are positive. For necessity, take a common
branch eigenvalue $\mu$. Branch eigenvectors have nonzero attaching
coordinates; rescale them so that the two attaching contributions
cancel. Their union, with zero at $\ell$, is a nonzero vector $z$ with
$Az=\mu z$. Let $S$ be the alternating diagonal sign matrix. The
zero-diagonal part $J=A-aE_{\ell\ell}$ satisfies $SJ=-JS$;
since $z_\ell=0$, also
$A(Sz)=-\mu Sz$. Thus the nonzero symmetric matrix
$X=zz^{\mathsf T}+(Sz)(Sz)^{\mathsf T}$ commutes with $A$.
Its entries at path edges and at $(\ell,\ell)$ vanish, so it
violates nSSP. This is the obstruction in \cite[Theorem 4.1]{C}.
For unit edges, the branch eigenvalues are
$2\cos(j\pi/\ell)$, $1\leq j<\ell$, and
$2\cos(k\pi/(n+1-\ell))$, $1\leq k<n+1-\ell$.
Since cosine is injective on $(0,\pi)$, a common value exists exactly
when $\gcd(\ell,n+1-\ell)>1$. An empty branch has no eigenvalues,
and the same criterion gives $\gcd(\ell,n+1)=1$ at either endpoint.
\end{proof}

\section{Loop enlargement, obstructions and consequences}
The nSSP condition is open as a condition on $A$. Indeed it says that
the linear map
\[
 F_A:X\longmapsto (A\circ X,\ AX^{\mathsf T}-X^{\mathsf T}A)
\]
is injective. In fixed bases this is full column rank, witnessed by a
nonzero maximal minor, and the entries of its matrix depend continuously
on $A$. Write $A_0$ for the integer one-loop witness and
$D=\sum_{j\in L\setminus\{\ell\}}E_{jj}$, where $E_{jj}$ has a single
unit diagonal entry. For every sufficiently large positive integer $N$,
$A_0+N^{-1}D$ has nSSP by openness. Since $F_{cA}=cF_A$ for $c\ne0$,
\[
 A_N=NA_0+D
\]
is a symmetric nSSP witness with exactly the required support and
positive integer nonzero entries.

If $n$ is even, $r+s=n-1$ is odd, so one branch is even for every
choice of loop. If $n$ is odd and $\ell$ is odd, both branches are even.
Thus the preceding construction proves the sufficiency assertions of
Theorem~\ref{thm:main}, including $n=1$ (the matrix $(1)$).

For necessity, if $L=\varnothing$, $X=I$ violates nSSP for every
matrix of the pattern. If $n=2k+1$ and all loops lie at even vertices,
order the odd vertices first and write an arbitrary matrix of the
pattern as
\[
 A=\begin{pmatrix}0&B\\ C&D\end{pmatrix},\qquad
 B\in\R^{(k+1)\times k},\ C\in\R^{k\times(k+1)}.
\]
There exist nonzero vectors $u,v\in\R^{k+1}$ with $Cu=0$ and
$v^{\mathsf T}B=0$. Extend them by zero on the even vertices.
Then $Au=0$ and $v^{\mathsf T}A=0$. The nonzero matrix
$X=vu^{\mathsf T}$ is supported on odd--odd positions, where $A$ is
zero, and $X^{\mathsf T}=uv^{\mathsf T}$ commutes with $A$.
This is the obstruction of \cite[Theorem 5.7(Cc)]{KO}, and completes
the proof of Theorem~\ref{thm:main}.

\begin{corollary}\label{cor:siblings}
For odd $n\geq3$, a double path with its sole loop at the central
vertex allows but does not require nSSP when $n\equiv1\pmod4$,
and does not allow nSSP when $n\equiv3\pmod4$.
For every $n\geq4$, the path with loops at all interior vertices
allows but does not require nSSP; for $1\leq n\leq3$ it does not allow nSSP.
The pattern with $n=1$ and its one loop requires nSSP.
Every directed superpattern of an allowing path also allows nSSP,
with a positive-integer witness.
\end{corollary}
\begin{proof}
The central vertex $(n+1)/2$ is odd exactly in the first congruence
class. In either odd class with $n\geq3$, the unit-edge matrix has
$\gcd((n+1)/2,n+1)>1$ and fails nSSP by
Corollary~\ref{cor:spectral}, proving nonrequirement.
For odd $n\geq5$, the interior contains the odd vertex $3$;
the even case needs only a nonempty loop set, whereas for $n=3$
the sole interior vertex is even. For $n=1,2$ the interior-loop set
is empty. To see nonrequirement for the
interior-loop family, let $C$ be the unit-edge path matrix with
unit loops at both endpoints, and let $E$ have every entry equal to $1$.
Every row and column sum of $C$ is $2$, so $CE=EC$.
Put $B=I-C$ and $Y=I-E+C$. Then $B$ has exactly the interior-loop
pattern, $Y\ne0$ for $n\geq3$, $B\circ Y=0$, and $BY=YB$.
This is also the construction behind \cite[Theorem 5.7(A)]{KO}.
The one-dimensional assertion follows directly from \eqref{eq:test}.
For a superpattern, let $D$ have value $1$ at every additional arc
and zero elsewhere. The same openness and scaling argument gives
$NA_0+D$ for all sufficiently large integers $N$.
\end{proof}
The last assertion is also a consequence of the stronger, similarity-
preserving Superpattern Lemma, \cite[Lemma 1.2]{KO}; see also
\cite[Theorem 5.4]{F} for its non-symmetric form.

\section{Source coverage and verification}
In the following table, ``proved'' refers to the written arguments
above; ``prior proof cited'' identifies results used with their original
scope. The allowing classification is also \cite[Theorem 6.1]{C}.
{\small
\begin{longtable}{@{}>{\raggedright\arraybackslash}p{.31\linewidth}>{\raggedright\arraybackslash}p{.64\linewidth}@{}}
\toprule
Source item in \cite{KO} & Coverage\\ \midrule
\endfirsthead
\toprule Source item in \cite{KO} & Coverage\\ \midrule
\endhead
Question 5.11, Section 5; journal p.~776 & Proved: every positive even
order and nonempty loop set (Theorem~\ref{thm:main}); prior proof
\cite[Corollary 6.2]{C}.\\
Problem 1.6, allowing & Proved: complete criterion in all orders
(Theorem~\ref{thm:main}); prior proof \cite[Theorem 6.1]{C}.\\
Problem 1.6, requiring & Not covered in general; central-loop and
interior-loop subfamilies are classified in Corollary~\ref{cor:siblings}.\\
Theorem 5.7(Cc) & Prior proof cited: odd-order obstruction; its
rank-one proof is included above.\\
Example 5.8; Theorem 5.7(A) & Proved: interior loops allow but do not
require nSSP for every $n\geq4$; $1\leq n\leq3$ does not allow it.\\
Example 5.9 & Proved: the unit-edge witness for $P_{6,\{3\}}$ has
nSSP because $\gcd(3,7)=1$; the example's signed failure is prior
proof cited.\\
Example 5.10; Theorem 5.7(Ca) & Proved: central-loop three-way
classification for every odd order; unit edges fail for every
odd $n\geq3$.\\
Example 5.12; Table 1 & Proved: their allowing boundary agrees
with Theorem~\ref{thm:main}; the remaining requiring distinctions
are prior proof cited.\\
Theorems 5.3, 5.5 and 5.7(B),(Cb); Corollary 5.6 & Prior proof
cited: the stated requiring families and co-central-loop nonrequirement
are retained with their original hypotheses.\\
Lemma 1.2; \cite[Theorem 5.4]{F} & Prior proof cited: the stronger
similarity-preserving Superpattern Lemma. The positive-integer
existence consequence is proved above.\\
Problem 1.8; Problem 2.9 & Not covered: the general minimum-arc
requiring questions concern all matrices, rather than existence
of path witnesses.\\
Reference [1] of \cite{KO}, Problem 6.5 & Prior proof cited: \cite[Example 1.4]{KO}
already gives the negative answer to the irreducible-tridiagonal
requiring question.\\
\bottomrule
\end{longtable}}
The written proof above establishes the full criterion and the
positive-integer and superpattern consequences. The pinned Lean file
\href{https://github.com/the-omega-institute/trureturing/blob/b260643d3d496ad9cd04bfc3557fc9b7499de976/D5/S3/Combinatorics/SpectralGraph/KoljancicOblakEvenPathNSSP.lean}{even-order Lean source}
contains the declaration \texttt{result}, whose proposition is precisely
the even-order existence statement, with unrestricted real test matrices.
The companion
\href{https://github.com/the-omega-institute/trureturing/blob/b260643d3d496ad9cd04bfc3557fc9b7499de976/D5/S3/Combinatorics/SpectralGraph/KoljancicOblakEvenPath.lean}{continuant Lean source}
defines that proposition as \texttt{claim} and provides
\texttt{branch\_variation\_injective},\allowbreak\ 
\texttt{joined\_variation\_injective},\allowbreak\ and
\texttt{joined\_path\_charpoly\_derivative} for the algebraic
recovery and characteristic-polynomial differential. More precisely,
the branch-recovery declaration applies to any branch length with
nonzero edge parameters. The joined-recovery declaration requires
\texttt{Odd (r + s)}, a nonzero loop and coprime continuants; it covers
the even-order case. The characteristic-polynomial derivative declaration
applies to arbitrary branch lengths and real edge and loop variations.
The proposition \texttt{result} asserts an exact-support real witness.
Symmetry, positive integrality, the odd-order extension, the unit-edge
criterion and the nonrequiring constructions are written proofs here.
The accompanying \texttt{verify.py} checks exact rational ranks for
all loop sets through order eight and the branch differential through
order ten. It also checks the unit-edge arithmetic criterion for all
78 single-loop positions through order twelve, and the explicit
interior-loop failure for orders three through twelve. These are
finite checks; the all-orders statements follow from the written proofs.

\paragraph{License.}
\href{https://creativecommons.org/licenses/by/4.0/}{Creative Commons
Attribution 4.0 International (CC BY 4.0).}

\begin{thebibliography}{9}
\bibitem{KO}
S. Koljan\v{c}i\'c and P. Oblak,
\emph{Combinatorial aspects of the non-symmetric strong spectral
property for graphs}, arXiv:2604.21552v1 (2026),
\href{https://arxiv.org/abs/2604.21552v1}{arXiv version 1}.
Journal version: \emph{Combinatorial aspects of the nonsymmetric strong
spectral property for graphs}, Electron. J. Linear Algebra \textbf{42}
(2026), 755--779,
\href{https://doi.org/10.13001/ela.2026.10523}{doi:10.13001/ela.2026.10523};
Question 5.11, p.~776.
\bibitem{F}
S. M. Fallat, H. Tracy Hall, J. C.-H. Lin and B. L. Shader,
\emph{The bifurcation lemma for strong properties in the inverse
eigenvalue problem of a graph}, Linear Algebra Appl. \textbf{648}
(2022), 70--87,
\href{https://doi.org/10.1016/j.laa.2022.04.009}{doi:10.1016/j.laa.2022.04.009}.
\href{https://arxiv.org/abs/2112.11712v2}{arXiv:2112.11712v2}.
\bibitem{C}
Carptopus,
\emph{Spectral separation and the non-symmetric strong spectral property
for looped double paths}, 19 August 2026, revised 22 August 2026,
Theorems 4.1, 5.2 and 6.1 and Corollary 6.2.
\href{https://raw.githubusercontent.com/Carptopus/MathExplore-Notes/e57ee007380cbc351eb371d8e991303cab5f1e48/research/nssp-looped-double-paths/nssp-looped-double-paths.tex}{Dated original source};
\href{https://doi.org/10.5281/zenodo.22045508}{doi:10.5281/zenodo.22045508};
Version 0.2, 22 August 2026,
\href{https://doi.org/10.5281/zenodo.22059007}{doi:10.5281/zenodo.22059007}.
\bibitem{MT}
J. Michor and G. Teschl,
\emph{Reconstructing Jacobi matrices from three spectra},
in J. Janas, P. Kurasov and S. Naboko (eds.),
\emph{Spectral Methods for Operators of Mathematical Physics},
Operator Theory: Advances and Applications \textbf{154},
Birkh\"auser, Basel (2004), 151--154,
\href{https://doi.org/10.1007/978-3-0348-7947-7_9}{doi:10.1007/978-3-0348-7947-7\_9}.
\href{https://arxiv.org/abs/math/0211051}{arXiv:math/0211051}.
\bibitem{VM}
K. N. Vander Meulen,
\emph{Sign patterns that require or allow the non-symmetric strong
spectral property}, abstract, joint work with M. Cavers and Z. Li,
27th Conference of the International Linear Algebra Society (2026), p.~71.
\href{https://ilas2026.math.vt.edu/docs/ILAS2026-Book-Of-Abstracts.pdf}{Conference abstracts}.
\end{thebibliography}
\end{document}
